\documentclass[10pt,a4paper]{amsart}
\usepackage{amsmath,amssymb,amsthm,graphicx}
\usepackage{verbatim}
\newenvironment{reference}{\comment}{\endcomment}
%\newenvironment{reference}{}{}
\newenvironment{slogan}{\comment}{\endcomment}
\newenvironment{history}{\comment}{\endcomment}

% For commutative diagrams we use Xy-pic
\usepackage[all]{xy}

% We use 2cell for 2-commutative diagrams.
\xyoption{2cell}
% Theorem environments.
%
\theoremstyle{plain}
\newtheorem{theorem}[subsection]{Theorem}
\newtheorem{proposition}[subsection]{Proposition}
\newtheorem{lemma}[subsection]{Lemma}

\theoremstyle{definition}
\newtheorem{definition}[subsection]{Definition}
\newtheorem{example}[subsection]{Example}
\newtheorem{exercise}[subsection]{Exercise}
\newtheorem{situation}[subsection]{Situation}

\theoremstyle{remark}
\newtheorem{remark}[subsection]{Remark}
\newtheorem{remarks}[subsection]{Remarks}

\numberwithin{equation}{subsection}

% Macros
%
\def\lim{\mathop{\mathrm{lim}}\nolimits}
\def\colim{\mathop{\mathrm{colim}}\nolimits}
\def\Spec{\mathop{\mathrm{Spec}}}
\def\Hom{\mathop{\mathrm{Hom}}\nolimits}
\def\Ext{\mathop{\mathrm{Ext}}\nolimits}
\def\SheafHom{\mathop{\mathcal{H}\!\mathit{om}}\nolimits}
\def\SheafExt{\mathop{\mathcal{E}\!\mathit{xt}}\nolimits}
\def\Sch{\mathit{Sch}}
\def\Mor{\mathop{\mathrm{Mor}}\nolimits}
\def\Ob{\mathop{\mathrm{Ob}}\nolimits}
\def\Sh{\mathop{\mathit{Sh}}\nolimits}
\def\NL{\mathop{N\!L}\nolimits}
\def\CH{\mathop{\mathrm{CH}}\nolimits}
\def\proetale{{pro\text{-}\acute{e}tale}}
\def\etale{{\acute{e}tale}}
\def\QCoh{\mathit{QCoh}}
\def\Ker{\mathop{\mathrm{Ker}}}
\def\Im{\mathop{\mathrm{Im}}}
\def\Coker{\mathop{\mathrm{Coker}}}
\def\Coim{\mathop{\mathrm{Coim}}}

% Boxtimes
%
\DeclareMathSymbol{\boxtimes}{\mathbin}{AMSa}{"02}

%
% Macros for moduli stacks/spaces
%
\def\QCohstack{\mathcal{QC}\!\mathit{oh}}
\def\Cohstack{\mathcal{C}\!\mathit{oh}}
\def\Spacesstack{\mathcal{S}\!\mathit{paces}}
\def\Quotfunctor{\mathrm{Quot}}
\def\Hilbfunctor{\mathrm{Hilb}}
\def\Curvesstack{\mathcal{C}\!\mathit{urves}}
\def\Polarizedstack{\mathcal{P}\!\mathit{olarized}}
\def\Complexesstack{\mathcal{C}\!\mathit{omplexes}}
% \Pic is the operator that assigns to X its picard group, usage \Pic(X)
% \Picardstack_{X/B} denotes the Picard stack of X over B
% \Picardfunctor_{X/B} denotes the Picard functor of X over B
\def\Pic{\mathop{\mathrm{Pic}}\nolimits}
\def\Picardstack{\mathcal{P}\!\mathit{ic}}
\def\Picardfunctor{\mathrm{Pic}}
\def\Deformationcategory{\mathcal{D}\!\mathit{ef}}
\usepackage{cleveref}
\crefname{originalreference}{original reference}{original references}
\Crefname{originalreference}{Original reference}{Original references}
\makeatletter
\newlabel{situation-quasi-compact-objects}{{0937}{source}{}{sourcepointer-1}{}}
\newlabel{situation-quasi-compact-objects@cref}{{[originalreference][0][]0937}{[1][0][]source}{}{}{}}
\makeatother
\crefname{originalreference}{original reference}{original references}
\Crefname{originalreference}{Original reference}{Original references}
\makeatletter
\newlabel{item-enough}{{0938}{source}{}{sourcepointer-2}{}}
\newlabel{item-enough@cref}{{[originalreference][0][]0938}{[1][0][]source}{}{}{}}
\makeatother
\crefname{originalreference}{original reference}{original references}
\Crefname{originalreference}{Original reference}{Original references}
\makeatletter
\newlabel{item-enough-qc}{{0939}{source}{}{sourcepointer-3}{}}
\newlabel{item-enough-qc@cref}{{[originalreference][0][]0939}{[1][0][]source}{}{}{}}
\makeatother
\crefname{originalreference}{original reference}{original references}
\Crefname{originalreference}{Original reference}{Original references}
\makeatletter
\newlabel{item-enough-qc-qs}{{093A}{source}{}{sourcepointer-4}{}}
\newlabel{item-enough-qc-qs@cref}{{[originalreference][0][]093A}{[1][0][]source}{}{}{}}
\makeatother
\crefname{originallemma}{original lemma}{original lemmas}
\Crefname{originallemma}{Original lemma}{Original lemmas}
\makeatletter
\newlabel{lemma-silly-quasi-compact}{{0G1W}{source}{}{sourcepointer-5}{}}
\newlabel{lemma-silly-quasi-compact@cref}{{[originallemma][0][]0G1W}{[1][0][]source}{}{}{}}
\makeatother
\title{Refining a presentation towards constructible sheaves on a site with enough quasi-compact objects}
\author{The Stacks Project Authors}
\date{stacks-project:sites-modules.tex:local-snapshot-20261012; GNU FDL 1.2 or later}
\begin{document}
\maketitle
\noindent\textit{Editorial note.} All delivered mathematical text is the author's own, copied verbatim from the pinned Stacks Project chapter file sites-modules.tex; no statement or proof step was rewritten, reordered or removed.\par\medskip
\noindent\textit{Editorial note.} The statement and the proof contain commutative diagrams typeset with Xy-pic; the retained two-slice preamble of this source supplies the native Xy-pic setup, so they compile as in the source.\par\medskip
\noindent\textit{Editorial note.} The counted span cites a situation and a refinement lemma, both recorded in assembly.external\_references with their canonical tags as navigation pointers.\par\medskip
\noindent\textit{Allowed standard prerequisites (editorial).}\begin{itemize}
\item Quasi-compact objects admit finite subcoverings of a basis covering. Stacks Project, Modules on Sites, Lemma lemma-silly-quasi-compact (tag 09B0). Provides the finite subfamilies K'\_i and M'\_j used at both refinement steps.
\item The quasi-compact objects situation and its enough-objects hypotheses. Stacks Project, Modules on Sites, Situation situation-quasi-compact-objects (tag 09AZ). Supplies the ambient hypotheses (items enough, enough-qc, enough-qc-qs) under which the refinements are possible.
\end{itemize}\medskip
\ifdefined\hypertarget\hypertarget{sourcepointer-1}{}\fi
\noindent\textit{Source reference.} Original reference 0937: sources/stacks-chow/tags.txt, source lines 11505--11505. Cross-chapter standard prerequisite cited by the author inside the counted span; it is used as a pointer and its own proof is not part of this item.\par\medskip
\ifdefined\hypertarget\hypertarget{sourcepointer-2}{}\fi
\noindent\textit{Source reference.} Original reference 0938: sources/stacks-chow/tags.txt, source lines 11506--11506. Cross-chapter standard prerequisite cited by the author inside the counted span; it is used as a pointer and its own proof is not part of this item.\par\medskip
\ifdefined\hypertarget\hypertarget{sourcepointer-3}{}\fi
\noindent\textit{Source reference.} Original reference 0939: sources/stacks-chow/tags.txt, source lines 11507--11507. Cross-chapter standard prerequisite cited by the author inside the counted span; it is used as a pointer and its own proof is not part of this item.\par\medskip
\ifdefined\hypertarget\hypertarget{sourcepointer-4}{}\fi
\noindent\textit{Source reference.} Original reference 093A: sources/stacks-chow/tags.txt, source lines 11508--11508. Cross-chapter standard prerequisite cited by the author inside the counted span; it is used as a pointer and its own proof is not part of this item.\par\medskip
\ifdefined\hypertarget\hypertarget{sourcepointer-5}{}\fi
\noindent\textit{Source reference.} Original lemma 0G1W: sources/stacks-chow/tags.txt, source lines 20036--20036. Cross-chapter standard prerequisite cited by the author inside the counted span; it is used as a pointer and its own proof is not part of this item.\par\medskip
\begin{lemma}
\label{lemma-change-presentation-towards-constructibles}
In Situation \ref{situation-quasi-compact-objects} assume
(\ref{item-enough}), (\ref{item-enough-qc}), and (\ref{item-enough-qc-qs})
hold. Let $\mathcal{O}$ be a sheaf of rings. Assume given a map
$$
\bigoplus\nolimits_{j = 1, \ldots, m} j_{V_j!}\mathcal{O}_{V_j}
\longrightarrow
\bigoplus\nolimits_{i = 1, \ldots, n} j_{U_i!}\mathcal{O}_{U_i}
$$
with $U_i$ and $V_j$ in $\mathcal{B}$, and coverings
$\{U_{ik} \to U_i\}_{k \in K_i}$ with $U_{ik} \in \mathcal{B}$.
Then there exist finite subsets $K'_i \subset K_i$ and
a finite set $L$ of $W_l \in \mathcal{B}$ and a commutative diagram
$$
\xymatrix{
\bigoplus_{l \in L} j_{W_l!}\mathcal{O}_{W_l} \ar[d] \ar[r] &
\bigoplus_{i = 1, \ldots, n} \bigoplus_{k \in K'_i}
j_{U_{ik}!}\mathcal{O}_{U_{ik}} \ar[d] \\
\bigoplus_{j = 1, \ldots, m} j_{V_j!}\mathcal{O}_{V_j} \ar[r] &
\bigoplus_{i = 1, \ldots, n} j_{U_i!}\mathcal{O}_{U_i}
}
$$
inducing an isomorphism on cokernels of the horizontal maps.
\end{lemma}

\begin{proof}
Since $U_i$ is quasi-compact, we may choose finite subsets $K'_i \subset K_i$
as in Lemma \ref{lemma-silly-quasi-compact}.
Then since
$\bigoplus_{i = 1, \ldots, n}
\bigoplus_{k \in K'_i} j_{U_{ik}!}\mathcal{O}_{U_{ik}} \to
\bigoplus_{i = 1, \ldots, n} j_{U_i!}\mathcal{O}_{U_i}$ is surjective,
we can find coverings $\{V_{jm} \to V_j\}_{m \in M_j}$ with
$V_{jm} \in \mathcal{B}$ such that we can find a commutative diagram
$$
\xymatrix{
\bigoplus_{j = 1, \ldots, m} \bigoplus_{m \in M_j}
j_{V_{jm}!}\mathcal{O}_{V_{jm}} \ar[d] \ar[r] &
\bigoplus_{i = 1, \ldots n} \bigoplus_{k \in K'_i}
j_{U_{ik}!}\mathcal{O}_{U_{ik}} \ar[d] \\
\bigoplus_{j = 1, \ldots, m} j_{V_j!}\mathcal{O}_{V_j} \ar[r] &
\bigoplus_{i = 1, \ldots, n} j_{U_i!}\mathcal{O}_{U_i}
}
$$
Since $V_j$ is quasi-compact, we can choose finite subsets
$M'_j \subset M_j$ as in Lemma \ref{lemma-silly-quasi-compact}.
Set
$$
L = \left(\coprod\nolimits_{i = 1, \ldots, n} K'_i \times K'_i \right)
\coprod
\left(\coprod\nolimits_{j = 1, \ldots, m} M'_j\right)
$$
and for $l = (k, k') \in K'_i \times K'_i \subset L$ set
$W_l = U_{ik} \times_{U_i} U_{ik'}$ and for
$l = m \in M'_j \subset L$ set $W_l = V_{jm}$.
Since we have the exact sequences of Lemma \ref{lemma-silly-quasi-compact}
for the families $\{U_{ik} \to U_i\}_{k \in K'_i}$
we conclude that we get a diagram as in the statement of the lemma
(details omitted), except that it is not yet clear that $W_l \in \mathcal{B}$.
However, since $W_l$ is quasi-compact for all $l \in L$
we do another application of Lemma \ref{lemma-silly-quasi-compact}
and find finite families of maps $\{W_{lt} \to W_l\}_{t \in T_l}$
with $W_{lt} \in \mathcal{B}$ such that
$\bigoplus j_{W_{lt}!}\mathcal{O}_{W_{lt}} \to j_{W_l!}\mathcal{O}_{W_l}$
is surjective. Then we replace $L$ by $\coprod_{l \in L} T_l$
and everything is clear.
\end{proof}

\end{document}
